\section*{第 \ref{ch:g6:plants-are-producers} 章——植物是生产者}

\begin{solution}{exo:g6:plants-are-producers:1}
\omterm{def:g6:exploring-our-environment:components}{矿物质}：非生命世界的物质——水、
岩石（还包括空气中的气体、\omterm{prop:g4:what-plants-need:needs}{无机盐}）。\omterm{def:g6:plants-are-producers:matter}{有机物}：
由\omterm{def:g6:exploring-our-environment:components}{生物}制造并构成其身体的物质——
木材、\omterm{def:g3:bones-and-muscles:muscle}{肌肉}（还包括糖、羊毛、脂肪）。
\end{solution}

\begin{solution}{exo:g6:plants-are-producers:2}
投入：来自土壤的水和\omterm{prop:g4:what-plants-need:needs}{无机盐}，来自空气的气体（二氧
化碳）。\omterm{prop:g3:balanced-diet:jobs}{能量}：光，由绿色物质捕获。产出：\omterm{def:g6:plants-are-producers:matter}{有机物}——
糖，然后是木材、叶、
根、\omterm{def:g2:from-seed-to-plant:seed}{种子}。车间：白天的绿叶。
\end{solution}

\begin{solution}{exo:g6:plants-are-producers:3}
\omterm{def:g6:exploring-our-environment:components}{矿物质}：露水、花岗岩、海盐。\omterm{def:g6:plants-are-producers:matter}{有机物}：蜂蜡、食用
油——而蜗牛壳是蜗牛自己制造的，由\omterm{def:g1:animals-around-us:animal}{动物}沉积下来的
\omterm{def:g6:plants-are-producers:matter}{无机物}质构成：由\omterm{def:g6:exploring-our-environment:components}{生物}体利用无机原料制成。
（两种答案，只要理由如此，
均给分。）
\end{solution}

\begin{solution}{exo:g6:plants-are-producers:4}
主要来自空气和水——从空气中吸收的气体
以及根吸收的水——在叶中借光的\omterm{prop:g3:balanced-diet:jobs}{能量}合成
；土壤只提供少量
\omterm{prop:g4:what-plants-need:needs}{无机盐}。
\end{solution}

\begin{solution}{exo:g6:plants-are-producers:5}
它吸收二氧化碳，释放氧气。凡是\omterm{def:g4:breathing:movements}{呼吸}的
\omterm{def:g6:exploring-our-environment:components}{生物}都呼出前者、需要后者——
包括我们自己。
\end{solution}

\begin{solution}{exo:g6:plants-are-producers:6}
马铃薯\omterm{ex:g4:plants-without-seeds:bulbtuber}{块茎}和郁金香\omterm{ex:g4:plants-without-seeds:bulbtuber}{鳞茎}（地下储备），
豆类肥厚的两片（\omterm{def:g2:from-seed-to-plant:seed}{种子}的食物储备），胡萝卜肥大的
根、富含糖分的甜菜——都是夏季制造并
储存的\omterm{def:g6:plants-are-producers:matter}{有机物}。
\end{solution}

\begin{solution}{exo:g6:plants-are-producers:7}
因为制造过程依靠光：没有光，
就没有\omterm{prop:g3:balanced-diet:jobs}{能量}来合成\omterm{def:g6:plants-are-producers:matter}{有机物}，无论土壤中有多少\omterm{def:g6:exploring-our-environment:components}{矿物质}等待。
光是唯一无法替代的投入——
\omterm{prop:g4:what-plants-need:needs}{无机盐}只是调味料，不是燃料。
\end{solution}

\begin{solution}{exo:g6:plants-are-producers:8}
因为\omterm{def:g3:food-chains:chain}{食物链}的第一个环节必须把新的\omterm{def:g6:plants-are-producers:matter}{有机物}
送入\omterm{def:g6:exploring-our-environment:components}{生物}界，而\omterm{def:g1:animals-around-us:animal}{动物}是\omterm{def:g5:food-webs:roles}{消费者}：它们只能
转化别人已经制成的\omterm{def:g6:plants-are-producers:matter}{有机物}。只有
\omterm{def:g5:food-webs:roles}{生产者}能用无机原料制造\omterm{def:g6:plants-are-producers:matter}{有机物}——所以
第一个环节总是绿色的。
\end{solution}

\begin{solution}{exo:g6:plants-are-producers:9}
面包：小麦——小麦的\omterm{def:g2:from-seed-to-plant:seed}{种子}；追溯路径的终点是
小麦叶。黄油：奶油——牛奶——奶牛——草；
终点是草叶。奶酪：牛奶——奶牛——
仍是草。三种食物，三（或两）片叶，都
离不开光。
\end{solution}

\begin{solution}{exo:g6:plants-are-producers:10}
干草是奶牛的食物——奶牛将其转化为
自身身体的\omterm{def:g6:plants-are-producers:matter}{有机物}。堆肥并不是\omterm{def:g1:plants-around-us:plant}{植物}的食物：它在腐烂时
为土壤补充\emph{无机}盐。\omterm{def:g1:plants-around-us:plant}{植物}随后
用这些\omterm{prop:g4:what-plants-need:needs}{无机盐}、空气、水和光
自行制造\omterm{def:g6:plants-are-producers:matter}{有机物}。奶牛吃现成的食物；\omterm{def:g1:plants-around-us:plant}{植物}自己制造食物。
\end{solution}

\begin{solution}{exo:g6:plants-are-producers:11}
有光时，\omterm{def:g1:plants-around-us:plant}{植物}白天的制造过程释放氧气，
吸收二氧化碳，与鱼的\omterm{def:g4:breathing:movements}{呼吸}相平衡——
气体交换是双向的，瓶内形成循环。黑暗中，
制造停止，而每个居住者（包括\omterm{def:g1:plants-around-us:plant}{植物}）
仍在\omterm{def:g4:breathing:movements}{呼吸}：氧气耗尽，二氧化碳积累，
这个密封的小世界便窒息了。
\end{solution}

\begin{solution}{exo:g6:plants-are-producers:12}
正确：树的物质并非来自土壤——他的
称重实验永远证明了这一点。不完整：水只是
其中一种主要成分；他无法知道从
空气中吸收的气体，它提供了木材的大部分物质。
现代补充：来自空气的二氧化碳（加上
土壤中少量的\omterm{prop:g4:what-plants-need:needs}{无机盐}）——水和空气共同作用，
在光的作用下合成。
\end{solution}

\begin{solution}{pb:g6:plants-are-producers:1}
\textbf{1.} 这些吨数是由
\omterm{def:g1:plants-around-us:plant}{植物}自己用无机原料制造的\omterm{def:g6:plants-are-producers:matter}{有机物}——水、
空气中的二氧化碳、土壤中的\omterm{prop:g4:what-plants-need:needs}{无机盐}——以光为
\omterm{prop:g3:balanced-diet:jobs}{能量}。袋装肥料补充\omterm{def:g6:exploring-our-environment:components}{矿物质}；空气和雨水
提供主要部分；\omterm{def:g1:plants-around-us:parts}{叶子}负责制造。

\textbf{2.} 马铃薯：地下储藏器官（\omterm{ex:g4:plants-without-seeds:bulbtuber}{块茎}）。
胡萝卜：储藏根。豆类\omterm{def:g2:from-seed-to-plant:seed}{种子}：\omterm{def:g2:from-seed-to-plant:seed}{种子}的食物储备，
为下一代\omterm{def:g1:plants-around-us:plant}{植物}储备好。生菜叶：作为工作
器官本身，即捕获光能的车间，可鲜食。

\textbf{3.} 真正的供能者是光：每克产量
都是白天在叶中借助光能积累形成的。

\textbf{4.} 就\omterm{prop:g4:what-plants-need:needs}{无机盐}而言：收获会带走它们，
若不加以补充，\omterm{def:g1:plants-around-us:plant}{植物}就会缺乏\omterm{def:g6:exploring-our-environment:components}{矿物质}，
好比沙地里的薄荷——用量虽小，
作用却不可或缺。

\textbf{5.} 暗柜实验（苍白的囚徒）：遮光布隔绝了光，
仅一周时间，植株便已显现——
发白、细长、制造能力弱。

\textbf{6.} \omterm{ex:g4:plants-without-seeds:bulbtuber}{块茎}由地上部分充实：叶片
白天制造糖类，多余部分被运送至下方，
储存于膨大的\omterm{ex:g4:plants-without-seeds:bulbtuber}{块茎}中。繁茂的功能叶意味着
盈余量大——地上部分好，\omterm{ex:g4:plants-without-seeds:bulbtuber}{块茎}就好。

\textbf{7.} 水。制造过程需要全部原料同时到位：
光提供\omterm{prop:g3:balanced-diet:jobs}{能量}，而非物质——根部若干涸，
即便天空晴朗，也无物可用以构建。

\textbf{8.} 翻入土中的根是归还土壤的\omterm{def:g6:plants-are-producers:matter}{有机物}：
腐烂后将释放其中的\omterm{def:g6:exploring-our-environment:components}{矿物质}，供后茬作物利用——
即
\cref{prop:g5:people-and-nature:lasting}中的归还规则，以\omterm{def:g1:plants-around-us:plant}{植物}物质偿还。
焚烧则是浪费。

\textbf{9.} 马铃薯植株 $\rightarrow$ 猪 $\rightarrow$ 一家
人。其全部\omterm{def:g6:plants-are-producers:matter}{有机物}最初都在马铃薯植株的叶中
制成——猪与一家人只是
对其进行转化。

\textbf{10.} 原料：来自土壤的水和\omterm{prop:g4:what-plants-need:needs}{无机盐}，
来自空气的二氧化碳——组装成叶片
物质。太阳的确切作用：提供\omterm{prop:g3:balanced-diet:jobs}{能量}，而非物质。她的
沙拉由水、空气和土壤构成，\emph{由阳光造就}——
诗意的赞誉当之无愧，原料清单却有误。

\textbf{11.} \omterm{def:g5:food-webs:roles}{分解者}——土壤中的\omterm{def:g6:exploring-our-environment:components}{生物}，从蚯蚓、潮虫
直至肉眼难见的微小\omterm{def:g6:exploring-our-environment:components}{生物}——以枯叶为食，
并将其\omterm{def:g6:plants-are-producers:matter}{有机物}逐步分解为\omterm{def:g6:exploring-our-environment:components}{矿物质}，
以补充土壤，供下一茬作物
取用。

\textbf{12.} 例如：“花园里的\omterm{def:g5:food-webs:roles}{生产者}构建\omterm{def:g6:plants-are-producers:matter}{有机物}
——即收获物——，原料为取自土壤、水和空气的\omterm{def:g6:exploring-our-environment:components}{矿物质}，
并以光作为\omterm{prop:g3:balanced-diet:jobs}{能量}。”
\end{solution}
